\section{CONCLUSION AND IMPLICATIONS}
This study compared static, daily, and real-time comfort aggregation under observed office occupancy dynamics.
RT-GCM achieved the highest predicted mean comfort probability across all tested rooms and group sizes.
SGCM improvement diminished to approximately zero as the group size grew, whereas DGCM and RT-GCM remained beneficial by excluding absent occupants.
The additional RT-GCM gain over DGCM was 6--8 percentage points, showing that within-day composition information can matter even after daily attendance is known.
Real-time information also implied operational consequences.
Setpoint changes shrank as daily mean occupancy increased but still exceeded a common 0.5$^\circ$C control step.
The actionable-update rate rose sharply with UTR but was insensitive to $K$.
Within this comfort-only simulation scope, the practical implication is selective deployment: static aggregation for stable membership, daily aggregation for variable attendance with moderate turnover, and real-time tracking where sparse, rapidly changing composition justifies the sensing and control effort.
UTR could serve as a screening metric for selecting the aggregation resolution.